\section{Shrinking tubes and limiting punctures}
\label{sec:hermitian-tube-puncture-separation}

Lengths and areas below use the real metrics associated to Hermitian metrics. The diameter of a punctured subset uses the global intrinsic length distance on the punctured manifold.

\begin{example}\label{ex:hermitian-tube-puncture-separation}
There are a compact complex surface $X$, a holomorphic map $\pi:X\to Y=\mathbb P^2$ contracting a curve $E$ to a point $p$, and smooth Hermitian metrics $g_j,g_*$ on $X$ with the following properties. The map is biholomorphic from $X\setminus E$ to $Y\setminus\{p\}$, and $g_j=g_*$ eventually on each compact subset of $X\setminus E$. Their associated Hermitian forms satisfy
\[
 \omega_j^2=\omega_*^2.
\]
The curve areas tend to zero, and the metrics have a uniform integrable radial bound near $E$. For a neighborhood basis $V_\delta$ of $p$, with $U_\delta=\pi^{-1}(V_\delta)$,
\[
 \lim_{\delta\downarrow0}\limsup_{j\to\infty}
       \operatorname{diam}^{\mathrm{int}}(U_\delta,g_j)=0.
\]
Nevertheless, for the intrinsic distance $d_0$ of $g_0=(\pi^{-1})^*g_*$ on $Y\setminus\{p\}$, there is $c>0$ such that
\[
 \operatorname{diam}_{d_0}(V_\delta\setminus\{p\})\geq c
\]
for every sufficiently small $\delta$.
\end{example}
\begin{proof}
Put $p=[1:0:0]$ and take the incidence surface
\[
 X=\bigl\{([z_0:z_1:z_2],[u:v])\in\mathbb P^2\times\mathbb P^1:
                     z_1v=z_2u\bigr\},
 \qquad \pi([z],[u:v])=[z].
\]
Off $p$, the second coordinate is $[z_1:z_2]$, so $\pi$ is biholomorphic there. Its exceptional curve is $E=\{p\}\times\mathbb P^1$. In the affine chart $z_0=1$, a chart of $X$ is
\[
 (t,w)\longmapsto([1:w:tw],[1:t]).
\]
The other chart is $([1:t'w':w'],[t':1])$, with $t'=1/t$ and $w'=tw$ on their overlap. These are the charts of the tautological line bundle $L=\mathcal O_{\mathbb P^1}(-1)$. They also show smoothness along $E$; smoothness elsewhere follows from the biholomorphism. Thus $X$ is a compact smooth complex surface. Its connected dense complement $X\setminus E$ makes it connected.

Let $q:L\to\mathbb P^1$ be the bundle projection. Give $L$ the norm induced by its tautological inclusion into $\mathbb P^1\times\mathbb C^2$. In the displayed chart its squared norm is
\[
 r^2=(1+|t|^2)|w|^2=|z_1|^2+|z_2|^2.
\]
Fix an area-one Hermitian metric $h$ on $\mathbb P^1$. The Chern connection of the fiber norm defines a smooth semipositive Hermitian tensor $k$ on $L$ by the squared norm of the covariant fiber differential. Explicitly, with $\lambda=1+|t|^2$,
\[
 k=\lambda\left|\mathrm{d}w+w\partial\log\lambda\right|^2.
\]
Under $t'=1/t$, $w'=tw$, the differential in this formula multiplies by $t$ and $\lambda$ divides by $|t|^2$, so the tensor is globally defined. It has complex rank one, is positive on nonzero vertical tangent vectors, and restricts to zero on $TE$. Consequently $q^*h+k$ is a smooth positive Hermitian metric on $L$.

Choose $0<\delta_0<1$ and extend this metric from $U_{\delta_0}=\{r<\delta_0\}$ to a smooth Hermitian metric $g_*$ on $X$. Such an extension follows by a smooth cutoff supported in a slightly larger bundle tube and a convex combination with any global Hermitian metric. It agrees exactly with $q^*h+k$ on $U_{\delta_0}$.

Choose $0<\varepsilon_j<1$ tending to zero, with $2\varepsilon_j^2<\delta_0$, and a smooth function $0\leq\chi\leq1$ on $[0,\infty)$, equal to zero on $[0,1]$ and to one on $[4,\infty)$. Set
\[
 f_j(r)=\varepsilon_j^2+(1-\varepsilon_j^2)
                     \chi(r^2/\varepsilon_j^4),
 \qquad
 g_j=f_jq^*h+f_j^{-1}k
\]
on $U_{\delta_0}$, and set $g_j=g_*$ elsewhere. The function $f_j$ is smooth at $E$, is at least $\varepsilon_j^2$, and equals one for $r\geq2\varepsilon_j^2$. Hence the metrics glue smoothly and are positive. They equal $g_*$ eventually on every compact subset of $X\setminus E$.

Let $\beta,\kappa$ be the Hermitian forms of $q^*h,k$. Their complex rank one gives $\beta^2=\kappa^2=0$. On the tube,
\[
 (f_j\beta+f_j^{-1}\kappa)^2
       =2\beta\wedge\kappa=(\beta+\kappa)^2.
\]
Outside it the metrics coincide. This proves the global volume identity. On $E$, the restriction is $g_j|_E=\varepsilon_j^2h$, so its area is $\varepsilon_j^2$.

For any unit fiber vector $a$, a radial ray $\gamma(s)=sa$ has $k$-squared speed one and zero base component. Thus its $g_j$-squared speed is $f_j(s)^{-1}$. If $s<2\varepsilon_j^2$, this is at most $\varepsilon_j^{-2}\leq2/s$; otherwise it is one, still at most $2/s$ for $s<\delta_0$. This is the uniform integrable radial bound with exponent $1/2$.

A radial segment of radius $r<\delta$ has length at most $r+2\varepsilon_j$: its portion below $2\varepsilon_j^2$ has speed at most $\varepsilon_j^{-1}$ and length at most $2\varepsilon_j$, and its remaining speed is one. Connect two points by their radial segments and an almost minimizing path on $E$. All three segments stay inside $U_\delta$, giving
\[
 \operatorname{diam}^{\mathrm{int}}(U_\delta,g_j)
 \leq2\delta+\varepsilon_j\bigl(4+\operatorname{diam}(\mathbb P^1,h)\bigr).
\]
Taking $j\to\infty$ before $\delta\downarrow0$ proves the asserted shrinking condition. The affine balls
\[
 V_\delta=\{[1:z_1:z_2]:|z_1|^2+|z_2|^2<\delta^2\}
\]
have precisely these tubes as their preimages.

Let $d_*$ be the global distance of the fixed smooth metric $g_*$. Choose distinct points $a,b\in E$, and put $c=d_*(a,b)>0$. Choose points $a_r,b_r$ of fiber norm $r>0$ over these base points. They tend to $a,b$ as $r\downarrow0$. Every path in $Y\setminus\{p\}$ lifts to a path in $X\setminus E$ with the same $g_0$-length as its $g_*$-length. Therefore
\[
 d_0(\pi(a_r),\pi(b_r))\geq d_*(a_r,b_r)\longrightarrow c.
\]
For any fixed $\delta>0$, these two images lie in $V_\delta$ when $r<\delta$. Taking the supremum and then $r\downarrow0$ gives its punctured diameter at least $c$.
\end{proof}
